% Two Blind Spots in Backtest Validation: Evidence from the U.S. Dollar Index
% Formatted with mdpiemu.sty, a visual emulation of the MDPI layout. For submission,
% replace \usepackage{mdpiemu} with the official MDPI class from the journal template
% (Definitions/mdpi.cls) and move the metadata into that class's macros.
\documentclass[10pt]{article}
\usepackage{mdpiemu}

\Title{Two Blind Spots in Backtest Validation: Evidence from the U.S. Dollar Index}
\Author{Aryan Meena $^{1}$ and Eugene Pinsky $^{1,*}$}
\AuthorNames{Aryan Meena and Eugene Pinsky}
\address{$^{1}$\quad Department of Computer Science, Metropolitan College, Boston University,
Boston, MA 02215, USA}
\corres{Correspondence: epinsky@bu.edu}
\abstract{Backtest validation is built almost entirely around one failure, that too many strategies
were tried. Using the U.S. Dollar Index and its constituent currency futures over 2010 to
2025, we document two failures this defence cannot detect. First, a widely used free data
source records session opening prices with a stale print defect that manufactures an
overnight anomaly of Sharpe $-1.36$ ($p<0.0001$) absent from exchange futures on the same
instrument over the same sessions. The standard diagnostics do not merely miss the artefact, they certify it. The
probability of backtest overfitting is 0.02 for the false result against 0.77 for the
correct data and before costs the artefact survives stepwise familywise error control,
because these tools discipline the selection of strategies and never inspect the inputs. An
exact binomial screen on the location of the open within the daily range detects the defect
at $p=3.7\times10^{-13}$ with no external reference series. Second, a regime conditioned
currency strategy reporting Sharpe 0.72 (Newey--West $t=2.78$) reverses once its full 216
configuration search is declared instead of the six reported. A correctly specified search
of 1,312 configurations drawn from published specifications then finds none exceeding the
Sharpe ratio the same search attains on block bootstrapped returns. Reporting that noise
ceiling would have exposed both failures at once.}
\keyword{backtest overfitting; data quality; deflated Sharpe ratio; foreign exchange; multiple testing; reproducibility; trading strategy evaluation}
\JEL{C12; C15; C52; F31; G14}

\begin{document}


\section{Introduction}

The modern defence of a quantitative backtest is a defence against one specific failure.
The analyst enumerates the strategies that were tried, then applies a correction for the
number of trials: White's Reality Check \citep{white2000}, Hansen's test of superior
predictive ability \citep{hansen2005}, the stepwise procedure of \citet{romanowolf2005},
the deflated Sharpe ratio of \citet{baileylopez2014}, the probability of backtest
overfitting of \citet{bailey2017}, or the haircut Sharpe ratio of \citet{harveyliu2015}.
This machinery is well developed and widely adopted.

It is also narrow. Every one of these tools takes the input data as given and asks a
question about the strategy set. Two failures therefore sit outside its field of view. The
input series may be wrong, in which case the correction is applied faithfully to a number
that never existed. Or the declared strategy set may be smaller than the set actually
explored, in which case the correction is applied to the wrong denominator.

This paper documents one of each, on the same instrument, with one codebase and then
shows what a correctly specified search on the same market produces. Our contribution is
not that vendor data contains errors. \citet{inceporter2006} established that for equity
returns two decades ago and their paper remains the standard reference. Our contribution
is what happens when that literature is placed next to the overfitting diagnostics that
postdate it, in particular the deflated Sharpe ratio and the probability of backtest
overfitting. The answer is that those diagnostics invert. On the corrupted series the probability of backtest
overfitting is 0.016 to 0.028, which reads as an exceptionally robust result. On the
correct series for the same instrument it is 0.765 to 0.779, which reads as pure noise
fitting. Both readings are internally correct. A stationary data defect keeps winning out
of sample for exactly the same reason a real anomaly does, so no overfitting diagnostic
can separate them.

Section~\ref{sec:data} describes the data. Section~\ref{sec:defect} isolates the defect
and shows what the diagnostics say about it, then gives a screen that detects it.
Section~\ref{sec:search} documents the second failure, an undeclared search space.
Section~\ref{sec:exhaustive} runs a correctly specified search and introduces the noise
ceiling. Section~\ref{sec:structure} asks whether anything in this market is profitable
once the analysis is driven by structure rather than by enumeration.
Section~\ref{sec:protocol} states the reporting protocol we recommend and
Section~\ref{sec:limits} the limitations.

\section{Data}\label{sec:data}

Table~\ref{tab:data} lists the four sources. The two U.S. Dollar Index series are the
centre of the paper. One is the ICE spot index as redistributed by a free provider, the
series on which the original result was found. The other is the ICE Dollar Index futures
contract obtained from a commercial vendor sourced from the exchange. They describe the
same underlying index over the same dates.

\begin{table}[H]
\caption{Data sources. The two dollar index series in the first two rows describe the same
underlying index over the same dates and are the comparison on which
Section~\ref{sec:defect} rests.}
\label{tab:data}
\begin{tabularx}{\textwidth}{lLlr}
\toprule
\textbf{Series} & \textbf{Source} & \textbf{Period} & \textbf{Sessions} \\
\midrule
Dollar index, free provider & Free redistribution of the ICE spot index & 2010-12 to 2025-12 & 3,932 \\
Dollar index futures & Commercial vendor, exchange sourced (DX) & 2010-12 to 2025-12 & 3,889 \\
Currency futures & CME contracts E6, B6, D6, J6, S6 & 2010-12 to 2025-12 & 3,800 \\
Spot rates and interest rates & Federal Reserve H.10 and OECD via FRED & 1999-01 to 2026-08 & 6,929 \\
\bottomrule
\end{tabularx}
\end{table}

The five currency futures are the CME contracts for the euro, sterling, Canadian dollar,
yen and Swiss franc. No liquid Swedish krona contract exists, so the tradeable universe is
five of the six index constituents. All five carry median daily volume above 25,000
contracts, contain no violation of the open, high, low, close ordering and are correctly
dated: their return correlation with the Federal Reserve noon fix peaks at lag zero.

Nearby continuous futures jump when the front contract changes. That jump is the spread
between two contracts and not a return any holder earns, so all returns are computed
within contract and the 61 roll sessions in each series are assigned a return of zero. This
also means that contract changes cannot contribute to the overnight leg examined in
Section~\ref{sec:defect}.

Every $t$ statistic reported below is corrected for heteroskedasticity and autocorrelation
following \citet{neweywest1987}, using a Bartlett kernel with the standard automatic
bandwidth $\lfloor 4(T/100)^{2/9}\rfloor$, which gives eight lags at this sample length.
The conclusions are insensitive to that choice: the overnight leg of the free series returns
$p<0.0001$ at eight, nine and twelve lags and the exchange series returns 0.328, 0.328 and
0.331.

\section{A Data Defect That the Diagnostics Certify}\label{sec:defect}

\subsection{The result and the first attempt to explain it}

Decomposing the daily return into an overnight leg from the previous close to the open and
a daytime leg from the open to the close produces, on the free series, an overnight leg
with a Sharpe ratio of $-1.36$ and $p<0.0001$. This is a large effect. On the exchange
futures for the same instrument the overnight leg has a Sharpe ratio of $+0.25$ with
$p=0.328$. Both figures are computed on the 3,793 overnight returns available in the 3,794
sessions for which both series report a complete bar, so the two are directly comparable.

A naive way to locate the source is to substitute one vendor's opening prices into the
other vendor's series. That test is invalid here. The two series carry a persistent level
offset of 0.041 index points, which is 6.7 per cent of a typical daily range, so a price
level substitution injects the offset into every return. Performing it produced a terminal
value of 1178 for one strategy, which is not a finding but an artefact of the offset.

\subsection{Isolating the open}

The valid version transfers the close to open \emph{return} rather than the price, so that
each row uses one vendor's price level throughout and the result is free of the offset.
Table~\ref{tab:swap} reports it.

\begin{table}[H]
\caption{Scale free decomposition. Each row uses a single vendor's price level throughout
and transfers only the close to open return, so the level offset between vendors cannot
contribute. The daytime leg is the residual, which makes the 24 hour return identical
across rows by construction. Terminal values start from 100 and the sample is the 3,793
overnight returns common to both series.}
\label{tab:swap}
\begin{tabularx}{\textwidth}{Lrrrrrr}
\toprule
 & \multicolumn{3}{c}{\textbf{Overnight leg}} & \multicolumn{3}{c}{\textbf{Daytime leg}} \\
\cmidrule(lr){2-4}\cmidrule(lr){5-7}
\textbf{Price level, opening return} & \textbf{Final} & \textbf{Sharpe} & \textbf{$p$} & \textbf{Final} & \textbf{Sharpe} & \textbf{$p$} \\
\midrule
Free, free open          & 61.68  & $-1.358$ & 0.0000 & 197.35 & 0.700 & 0.0086 \\
Free, exchange open      & 108.08 & $+0.252$ & 0.3280 & 112.62 & 0.151 & 0.5534 \\
Exchange, exchange open  & 108.08 & $+0.252$ & 0.3280 & 112.22 & 0.146 & 0.5668 \\
Exchange, free open      & 61.68  & $-1.358$ & 0.0000 & 196.65 & 0.678 & 0.0103 \\
\bottomrule
\end{tabularx}
\end{table}

The overnight leg follows whichever vendor's open is used and is unchanged by whose closes
are used. This is as close to a controlled experiment as observational data allows. The
opening prices are the entire source of the effect.

Table~\ref{tab:openloc} shows what is wrong with them. With the level offset removed, the
free series' open lies outside the exchange's true high to low range on 7.30 per cent of
sessions. On the 204 sessions where the free open equals its own daily low it sits below
the exchange low, at $-0.073$ of the exchange range and falls outside that range on 56.9
per cent of those sessions against 3.5 per cent on ordinary sessions. A price outside the
session range is not a price at which the instrument traded. The free series also records
an open exactly equal to the previous close on 6.84 per cent of sessions against 1.34 per
cent for the exchange series, which is the signature of a stale print.

\begin{table}[H]
\caption{Location of the recorded open within the exchange session range, with the level
offset removed. A value of 0 is the exchange low and 1 the exchange high. The sessions on
which the free open equals its own extreme are the sessions that carry the anomaly.}
\label{tab:openloc}
\begin{tabularx}{\textwidth}{Lrrr}
\toprule
\textbf{Session group} & \textbf{$n$} & \textbf{Open location} & \textbf{Outside range} \\
\midrule
Free open at its own low  & 204   & $-0.073$ & 56.9\% \\
Free open at its own high & 82    & $+0.994$ & 45.1\% \\
All other sessions        & 3,508 & $+0.494$ & 3.5\% \\
\midrule
Exchange open, all sessions & 3,889 & $+0.498$ & 0.00\% \\
\bottomrule
\end{tabularx}
\end{table}

\subsection{What the diagnostics say}

The 24 strategy grid of the original study was rerun on both series. On exchange data no
strategy beats buy and hold at any cost level, the analysis of variance across the 24
gives $F(23,93312)=0.306$ with $p=0.9994$ and the largest Sharpe ratio in the grid,
$+0.269$, is below the $+0.316$ that the null expects the maximum of 24 trials to produce.
The walk forward over 13 out of sample years ends at 97.69 against buy and hold at 122.64.
There is nothing in the exchange data.

Table~\ref{tab:diagnostics} is the central result of this section. Applied to the corrupted
series at zero cost, Hansen's test rejects at $p=0.013$, three configurations survive the
stepwise procedure of \citet{romanowolf2005} and the probability of backtest overfitting
is 0.016. Applied to the correct series the same diagnostics report 1.000, zero survivors
and 0.779. The diagnostics rank the false result far above the true one.

\begin{table}[H]
\caption{Multiple testing and overfitting diagnostics applied to the same 24 strategy grid
on the two dollar index series, at zero transaction cost. RW is the number of
configurations surviving the stepwise familywise error procedure of
\citet{romanowolf2005}. PBO is the probability of backtest overfitting from the
combinatorially symmetric cross validation of \citet{bailey2017}, reported at 12 and 16
blocks. Lower PBO indicates a result that persists out of sample.}
\label{tab:diagnostics}
\begin{tabularx}{\textwidth}{Lrrrr}
\toprule
\textbf{Series} & \textbf{White RC $p$} & \textbf{Hansen SPA $p$} & \textbf{RW survivors} & \textbf{PBO} \\
\midrule
Free provider (the artefact) & 0.1032 & 0.0132 & 3 of 24 & 0.016 to 0.028 \\
Exchange futures (correct)   & 0.9924 & 1.0000 & 0 of 24 & 0.765 to 0.779 \\
\bottomrule
\end{tabularx}
\end{table}

Both readings are correct on their own terms. The probability of backtest overfitting asks
whether a strategy selected in sample continues to win out of sample. A persistent data
defect continues to win, because it is present in both halves of the sample. To every
overfitting diagnostic a stationary artefact and a real anomaly are the same object. This
is a structural limitation of the entire family of tools and not a defect in any one of
them.

\subsection{A screen that requires no reference data}

The comparison above required a second, trustworthy series. Most analysts do not have one.
We therefore propose a test that uses only the series being examined.

Under the null that a recorded open is an unbiased draw from the prices visited during the
session, the open is equally likely to coincide with the session low as with the session
high. Counting the two events gives an exact binomial test with no free parameters:
\begin{equation}
n_{\text{low}} \sim \mathrm{Binomial}\!\left(n_{\text{low}}+n_{\text{high}},\ \tfrac{1}{2}\right).
\label{eq:binom}
\end{equation}

Table~\ref{tab:screen} applies Equation~\eqref{eq:binom} to eight series. It clears the
exchange series at $p=0.788$ and flags the corrupted one at $p=3.7\times10^{-13}$.

\begin{table}[H]
\caption{The open location binomial screen of Equation~\eqref{eq:binom} applied to eight
series. Mean position is the average location of the open within the session range, where
0 is the low and 1 the high. The screen clears the series known to be sound and rejects the
series known to be corrupt, using only that series' own open, high and low.}
\label{tab:screen}
\begin{tabularx}{\textwidth}{Lrrrr}
\toprule
\textbf{Series} & \textbf{At low} & \textbf{At high} & \textbf{Binomial $p$} & \textbf{Mean position} \\
\midrule
Dollar index futures, exchange & 64  & 60  & 0.788 & 0.499 \\
Dollar index, free provider    & 204 & 82  & $3.68\times10^{-13}$ & 0.479 \\
EURUSD, free provider          & 90  & 56  & 0.0061 & 0.501 \\
GBPUSD, free provider          & 69  & 58  & 0.375  & 0.504 \\
JPYUSD, free provider          & 125 & 103 & 0.164  & 0.489 \\
CADUSD, free provider          & 66  & 52  & 0.231  & 0.506 \\
CHFUSD, free provider          & 62  & 44  & 0.098  & 0.505 \\
SEKUSD, free provider          & 44  & 45  & 1.000  & 0.508 \\
\bottomrule
\end{tabularx}
\end{table}

Two honest qualifications. The screen also flags EURUSD at $p=0.0061$, which barely
survives a Bonferroni correction across the eight tests at a threshold of 0.00625. We
independently established that these spot series are date stamped one trading day late,
which is consistent with a true positive, but we cannot prove it is not a multiplicity
artefact. And the screen is validated here against exactly one instrument with a known
good reference. It is a promising diagnostic and not an established one.

We note also that the mean position of the open within the range, at 0.479 for the
corrupted series against 0.499 for the exchange series, is a weak discriminator. The count
of extremes carries the information, not the average location.

\section{An Undeclared Search Space}\label{sec:search}

The second failure needs no bad data. A regime conditioned cross sectional currency
strategy, in which six currencies are ranked each quarter and the two weakest are held with
a sign set by the 126 day momentum of the dollar index, reports a Sharpe ratio of 0.72 with
a Newey--West $t$ statistic of 2.78 and $p=0.0055$ on Federal Reserve spot data.

Three defects in the committed implementation had to be corrected first. The pandas call
\texttt{to\_period('2Q')} silently returns quarterly periods, so the reported semiannual
row was a byte identical duplicate of the quarterly row. The grid loop was wrapped in a
bare exception handler, which concealed that the annual frequency alias had been renamed
between pandas versions, so every annual configuration was silently discarded. The
construction \texttt{np.where(mom < 0, 1, -1)} maps a missing value to $-1$, so the first
126 sessions, in which no momentum can yet be computed, were traded as though the dollar
were rising; correcting this alone moves the headline Sharpe ratio from 0.54 to 0.46. And
the constituent price files are date stamped one trading day late relative to the index file.

Table~\ref{tab:declared} reports what happens when the search space is declared in full.
The original account reported six configurations. Frequency, cross sectional slice,
cardinality, the regime filter and the momentum lookback are all free parameters and
enumerating them gives 144 configurations on the tradeable five currency panel and 216 on
the six currency panels.

\begin{table}[H]
\caption{Multiple testing over the declared search space. The strategy is a self financing
long and short book, so the correct null is zero excess return; the committed code tested
against a long dollar position, which is shown for comparison. The number after each
universe is the count of currencies in it. No configuration survives the stepwise procedure
in any universe once the full search is declared.}
\label{tab:declared}
\begin{tabularx}{\textwidth}{Lrrrrrrr}
\toprule
 & & & \multicolumn{2}{c}{\textbf{Null: zero}} & \multicolumn{2}{c}{\textbf{Null: long dollar}} & \\
\cmidrule(lr){4-5}\cmidrule(lr){6-7}
\textbf{Universe} & \textbf{Configs} & \textbf{Sharpe} & \textbf{RC $p$} & \textbf{SPA $p$} & \textbf{RC $p$} & \textbf{SPA $p$} & \textbf{DSR} \\
\midrule
Exchange futures, 5 & 144 & 0.458 & 0.2333 & 0.3833 & 0.6730 & 0.7647 & 0.227 \\
Futures plus krona, 6 & 216 & 0.498 & 0.2203 & 0.2830 & 0.4373 & 0.5933 & 0.217 \\
Federal Reserve spot, 6 & 216 & 0.722 & 0.0963 & 0.1570 & 0.3173 & 0.4727 & 0.613 \\
\bottomrule
\end{tabularx}
\end{table}

The deflated Sharpe ratio is a confidence level and must exceed 0.95. It reaches 0.613 at
best. The expected maximum Sharpe ratio under the null across 144 configurations is
$+0.654$, above the $+0.458$ actually observed on tradeable instruments. No universe now
clears a conventional threshold on any test.

The configuration is also a spike rather than a plateau. Holding everything else at the
reported values on the tradeable five currency panel and moving one parameter at a time,
the regime lookback gives $-0.319$ with no overlay, $+0.136$ at 63 days, $+0.458$ at 126
days and $+0.259$ at 252 days. The rebalance frequency gives $-0.234$ daily, $-0.081$
weekly, $+0.047$ monthly, $+0.458$ quarterly, $+0.112$ semiannually and $+0.456$ annually.
Only two of six frequencies are materially positive and the neighbours of the reported
lookback lose most of the effect, which describes a noisy surface rather than a mechanism.

Finally, an ensemble that holds every configuration in a family at equal weight selects
nothing and therefore requires no correction. Averaging the net returns of the 36 regime
signed configurations in the bottom slice, which charges each member its own full turnover
and so is conservative, gives a Sharpe ratio of 0.148 with $p=0.57$. Widening the ensemble
to all 108 regime signed configurations gives $-0.160$ with $p=0.53$.

\section{A Correctly Specified Search}\label{sec:exhaustive}

\subsection{Validating the harness before believing it}

Before running any search we validated the evaluation code itself, because a paper about
invalid backtests should demonstrate that its own backtest is valid.
Table~\ref{tab:harness} lists the ten checks. The most important is the first. Weights
formed from the same day's realised return, applied without the one period shift, must
produce an absurd Sharpe ratio; if they do not, the shift is broken and nothing downstream
means anything.

\begin{table}[H]
\caption{Validation of the evaluation harness. All ten checks pass; the eight rows below
group the four turnover and statistic checks into two lines. The lookahead injection is the
critical one, because the harness must be demonstrably capable of detecting a strategy that
cheats.}
\label{tab:harness}
\begin{tabularx}{\textwidth}{Lr}
\toprule
\textbf{Check} & \textbf{Result} \\
\midrule
Perfect foresight weights applied without the shift & Sharpe $+24.34$ \\
The same weights applied through the harness        & Sharpe $-0.11$ \\
One extra lag applied to 14 real families           & none collapses \\
Full search on simulated zero drift returns, RC $p$ & 0.20, 0.07, 0.69, 0.38 \\
Observed best of $N$ on noise against theory        & 0.823 against 0.846 \\
Sharpe, Newey--West $t$, PSR, cost arithmetic       & reproduced to machine precision \\
Buy and hold turnover; daily sign flipper turnover  & 0.00; exactly 2.0 per session \\
Winning configuration rebuilt from first principles & maximum difference $0.0\mathrm{e}{+}00$ \\
\bottomrule
\end{tabularx}
\end{table}

\subsection{The search and the noise ceiling}

We evaluated 1,312 configurations built from published specifications: time series momentum
\citep{moskowitz2012}, cross sectional currency momentum \citep{menkhoff2012}, carry
\citep{lustig2011,koijen2018}, value \citep{asness2013}, moving average rules
\citep{brock1992}, channel breakouts, open interest growth \citep{hongyogo2012} and
volatility scaling \citep{moreiramuir2017}, each crossed with four rebalance frequencies,
the regime overlay at three lookbacks and with or without volatility targeting. Costs are
charged at one basis point per unit of turnover.

The natural benchmark is not zero. It is the Sharpe ratio this same search attains on
returns that contain no predictable structure. We construct that benchmark by resampling
the actual returns in blocks using the stationary bootstrap of
\citet{politisromano1994}, which preserves the marginal distributions, the fat tails and
the cross sectional correlation while destroying the time series structure a signal could
exploit, then running the entire search on the resampled panel. We call the resulting
maximum the \emph{noise ceiling}. It is not a new statistic. It is a Monte Carlo reality
check expressed in the units a practitioner reads.

\begin{table}[H]
\caption{The correctly specified search against its noise ceiling. The ceiling is the best
Sharpe ratio the identical 1,312 configuration search attains on block bootstrapped
returns, averaged over six resamples. No configuration on the real data exceeds it and
every correction agrees.}
\label{tab:ceiling}
\begin{tabularx}{\textwidth}{Lr}
\toprule
\textbf{Quantity} & \textbf{Value} \\
\midrule
Best configuration on the real data, net of costs & $+0.653$ \\
Noise ceiling, Gaussian null                      & $+0.752$ \\
Noise ceiling, block bootstrap of the real returns& $+0.712$ \\
Expected best of 1,312 under the null             & $+0.865$ \\
Configurations exceeding the ceiling              & \textbf{0 of 1,312} \\
\midrule
White Reality Check $p$                           & 0.659 \\
Hansen SPA $p$                                    & 0.773 \\
Romano--Wolf survivors                            & 0 \\
Deflated Sharpe ratio (needs $>0.95$)             & 0.203 \\
Harvey--Liu adjusted $p$, Bonferroni and BHY      & 1.000, 1.000 \\
Probability of backtest overfitting               & 0.556 \\
\bottomrule
\end{tabularx}
\end{table}

Table~\ref{tab:ceiling} gives the result. The best configuration found in fifteen years of
exchange data is \emph{below} what the same search finds in random data. The holdout is
consistent: selecting on 2011 to 2018 and evaluating once on 2019 to 2025, the in sample
winner ranks 163rd of 1,312 out of sample and the top ten in sample configurations average a
Sharpe ratio of $+0.019$.

Family ensembles, which select nothing, are reported in Table~\ref{tab:families}. Not one
is significant in either direction.

\begin{table}[H]
\caption{Family ensembles, equal weight within each family, costed in weight space so that
trades offsetting across members are not charged. Selecting nothing means no correction is
required, so these are the fair test of whether a family of ideas has value.}
\label{tab:families}
\begin{tabularx}{\textwidth}{Lrr@{\hskip 2em}Lrr}
\toprule
\textbf{Family} & \textbf{Sharpe} & \textbf{$p$} & \textbf{Family} & \textbf{Sharpe} & \textbf{$p$} \\
\midrule
Cross sectional momentum & $+0.218$ & 0.37 & Time series momentum & $-0.214$ & 0.37 \\
Value                    & $+0.204$ & 0.43 & Volume               & $-0.232$ & 0.37 \\
Carry                    & $+0.135$ & 0.59 & Breakout             & $-0.289$ & 0.23 \\
Reversal                 & $+0.133$ & 0.58 & Moving average       & $-0.341$ & 0.16 \\
Open interest            & $+0.094$ & 0.72 & Cross sectional reversal & $-0.360$ & 0.14 \\
\midrule
\multicolumn{3}{l}{All 1,312 configurations} & & $-0.176$ & 0.48 \\
\bottomrule
\end{tabularx}
\end{table}

\section{Is Anything Here Profitable?}\label{sec:structure}

An exhaustive search that finds nothing invites the objection that the search was the wrong
instrument. We therefore also proceeded analytically. The structure of the data is
unambiguous. The first order autocorrelation of returns is between $-0.004$ and $+0.033$,
variance ratios lie between 0.83 and 1.02 and an autoregression of returns on their own
lag gives $R^2$ between 0.0000 and 0.0011. The same autoregression on log realised
volatility gives $R^2$ between 0.966 and 0.972. Volatility is roughly three orders of
magnitude more forecastable than direction.

If direction cannot be forecast then only a genuine risk premium or a relationship true by
construction can pay. We tested both. The dollar index is an exact deterministic function
of its six constituents and the index future shares an expiry code with the five currency
futures on 98.2 per cent of sessions, so a matched synthetic basket is a no arbitrage
relationship. It is also efficiently priced: the convergence trade returns a Sharpe ratio
of $+0.09$ gross and $-0.20$ at one basis point per leg, because six legs are required
against a spread whose daily change has a standard deviation of five basis points.

Carry is the more serious candidate. The roll gap between the expiring and deferred
contract is the market's own quoted interest differential, which we validated against
Federal Reserve interbank differentials to within 0.43 percentage points per currency. On
the fifteen year futures panel it returns a Sharpe ratio of 0.377 and the mechanism
verifies: of the 1.34 per cent annual return, 1.10 per cent is the differential actually
held and 0.23 per cent is spot movement, which is the forward premium puzzle behaving as
described. Table~\ref{tab:carry} shows why we do not claim it.

\begin{table}[H]
\caption{Term structure carry on two samples. The futures panel begins in December 2010,
immediately after the 2008 carry drawdown. Extending the sample to April 2002 using
Federal Reserve spot and OECD interbank rates for all six constituents reverses the sign.}
\label{tab:carry}
\begin{tabularx}{\textwidth}{Lrr}
\toprule
\textbf{Sample} & \textbf{Currencies} & \textbf{Sharpe} \\
\midrule
Futures panel, 2011 to 2025            & 5 & $+0.377$ \\
Federal Reserve spot, 2002 to 2026     & 6 & $-0.249$ \\
\midrule
\multicolumn{3}{l}{\textit{Sub periods of the long sample}} \\
2002 to 2007 & 6 & $+0.418$ \\
2008 to 2013 & 6 & $-0.760$ \\
2014 to 2019 & 6 & $-0.246$ \\
2020 to 2026 & 6 & $-0.022$ \\
\bottomrule
\end{tabularx}
\end{table}

The canonical factor set on the long sample gives carry $-0.249$, momentum $+0.028$ and
value $+0.096$, with an equal risk combination of $-0.085$, against a noise ceiling whose
ninetieth percentile is $+0.41$. Volatility targeting, the one quantity that is
forecastable, changes the Sharpe ratio by $-0.009$ on average across seven series over 27
years and improves three of them, which is inside the bootstrap null. What it does deliver
is a reduction in the volatility of realised volatility from 0.025 to 0.015. That is risk
control and not return generation.

We report no profitable strategy. The reason is structural rather than a failure of effort.
There is no directional predictability, the only candidate risk premium in this universe is
negative over the full available sample and five to six highly correlated currencies with a
common factor explaining 54.8 per cent of variance cannot support a cross sectional
strategy.

\section{A Reporting Protocol}\label{sec:protocol}

Three additions, none expensive.

\textbf{Screen the inputs.} Report the binomial statistic of Equation~\eqref{eq:binom}, the
frequency with which the open equals the previous close and the count of sessions
violating the open, high, low, close ordering. These cost one pass over the data and would
have stopped the first failure before any strategy was run.

\textbf{Declare the whole search.} Report every parameter that was varied and the total
count of configurations evaluated, including those discarded. Report the ensemble across
each family alongside the selected configuration, since an ensemble selects nothing and
needs no correction.

\textbf{Report the noise ceiling.} Run the complete search once on block bootstrapped
returns and report the best Sharpe ratio it produces. Any result below that number is
indistinguishable from luck at that search size. In this study the ceiling was 0.71 to 0.75
for the 1,312 configuration search, 0.38 for a five strategy carry family and 0.41 for the
canonical factor set. Every real result fell below its own ceiling.

\section{Limitations}\label{sec:limits}

The evidence is one currency complex and two vendors. Nothing here generalises to free data
as a category or to this vendor across other instruments. The mechanism behind the defect
is inferred rather than observed: a stale print, a first electronic tick and a misfielded
session extreme are all consistent with what we see and we cannot inspect the vendor's
pipeline. The commercial series is itself a redistributor rather than the exchange and its
closing field is a last price whose settlement status we could not confirm. The screen is
validated against a single instrument. The noise ceiling is a presentation of an existing
idea rather than a new statistic. Finally, the absence of a profitable strategy is a
statement about six daily G10 series over fifteen to twenty seven years and not about
foreign exchange markets in general; emerging market carry, options and intraday data all
lie outside what we can test.

\section{Conclusions}

Two failures sit outside the reach of the standard defence of a backtest. A corrupted input
is not merely invisible to the overfitting diagnostics, it is actively certified by them,
because a stationary data defect persists out of sample exactly as a real anomaly does. An
undeclared search space causes the correction to be applied to the wrong denominator and a
strategy reporting $t=3.09$ can fail every test once the full search is stated. Both
failures are cheap to detect. A one line binomial test on the location of the open catches
the first and a single run of the complete search on block bootstrapped returns catches the
second, because it tells the analyst what their own search produces on data that contains
nothing.

\authorcontributions{Conceptualization, A.M. and E.P.; methodology, A.M.; software, A.M.;
validation, A.M.; writing, A.M. and E.P.; supervision, E.P. All authors have read and agreed
to the published version of the manuscript.}
\funding{This research received no external funding.}
\dataavailability{All code and the derived data required to reproduce every number in this
paper will be deposited in a public repository on acceptance and are available to reviewers
and to the editor on request in the interim. Federal Reserve H.10 and OECD interbank series
are publicly available through FRED. The exchange sourced futures data are commercial and
are subject to the vendor's redistribution terms.}
\conflictsofinterest{The authors declare no conflict of interest.}

\begin{thebibliography}{99}
\bibitem[Asness et al.(2013)]{asness2013} Asness, C.S.; Moskowitz, T.J.; Pedersen, L.H.
Value and Momentum Everywhere. \emph{J. Finance} \textbf{2013}, \emph{68}, 929--985.

\bibitem[Bailey and Lopez de Prado(2014)]{baileylopez2014} Bailey, D.H.; Lopez de Prado, M.
The Deflated Sharpe Ratio: Correcting for Selection Bias, Backtest Overfitting and
Non-Normality. \emph{J. Portf. Manag.} \textbf{2014}, \emph{40}, 94--107.

\bibitem[Bailey et al.(2017)]{bailey2017} Bailey, D.H.; Borwein, J.; Lopez de Prado, M.;
Zhu, Q.J. The Probability of Backtest Overfitting. \emph{J. Comput. Finance} \textbf{2017},
\emph{20}, 39--69.

\bibitem[Brock et al.(1992)]{brock1992} Brock, W.; Lakonishok, J.; LeBaron, B. Simple
Technical Trading Rules and the Stochastic Properties of Stock Returns. \emph{J. Finance}
\textbf{1992}, \emph{47}, 1731--1764.

\bibitem[Hansen(2005)]{hansen2005} Hansen, P.R. A Test for Superior Predictive Ability.
\emph{J. Bus. Econ. Stat.} \textbf{2005}, \emph{23}, 365--380.

\bibitem[Harvey and Liu(2015)]{harveyliu2015} Harvey, C.R.; Liu, Y. Backtesting.
\emph{J. Portf. Manag.} \textbf{2015}, \emph{42}, 13--28.

\bibitem[Hong and Yogo(2012)]{hongyogo2012} Hong, H.; Yogo, M. What Does Futures Market
Interest Tell Us about the Macroeconomy and Asset Prices? \emph{J. Financ. Econ.}
\textbf{2012}, \emph{105}, 473--490.

\bibitem[Ince and Porter(2006)]{inceporter2006} Ince, O.S.; Porter, R.B. Individual Equity
Return Data from Thomson Datastream: Handle with Care! \emph{J. Financ. Res.} \textbf{2006},
\emph{29}, 463--479.

\bibitem[Koijen et al.(2018)]{koijen2018} Koijen, R.S.J.; Moskowitz, T.J.; Pedersen, L.H.;
Vrugt, E.B. Carry. \emph{J. Financ. Econ.} \textbf{2018}, \emph{127}, 197--225.

\bibitem[Lustig et al.(2011)]{lustig2011} Lustig, H.; Roussanov, N.; Verdelhan, A. Common
Risk Factors in Currency Markets. \emph{Rev. Financ. Stud.} \textbf{2011}, \emph{24},
3731--3777.

\bibitem[Menkhoff et al.(2012)]{menkhoff2012} Menkhoff, L.; Sarno, L.; Schmeling, M.;
Schrimpf, A. Currency Momentum Strategies. \emph{J. Financ. Econ.} \textbf{2012},
\emph{106}, 660--684.

\bibitem[Moreira and Muir(2017)]{moreiramuir2017} Moreira, A.; Muir, T.
Volatility-Managed Portfolios. \emph{J. Finance} \textbf{2017}, \emph{72}, 1611--1644.

\bibitem[Moskowitz et al.(2012)]{moskowitz2012} Moskowitz, T.J.; Ooi, Y.H.; Pedersen, L.H.
Time Series Momentum. \emph{J. Financ. Econ.} \textbf{2012}, \emph{104}, 228--250.

\bibitem[Newey and West(1987)]{neweywest1987} Newey, W.K.; West, K.D. A Simple, Positive
Semi-Definite, Heteroskedasticity and Autocorrelation Consistent Covariance Matrix.
\emph{Econometrica} \textbf{1987}, \emph{55}, 703--708.

\bibitem[Politis and Romano(1994)]{politisromano1994} Politis, D.N.; Romano, J.P. The
Stationary Bootstrap. \emph{J. Am. Stat. Assoc.} \textbf{1994}, \emph{89}, 1303--1313.

\bibitem[Romano and Wolf(2005)]{romanowolf2005} Romano, J.P.; Wolf, M. Stepwise Multiple
Testing as Formalized Data Snooping. \emph{Econometrica} \textbf{2005}, \emph{73},
1237--1282.

\bibitem[White(2000)]{white2000} White, H. A Reality Check for Data Snooping.
\emph{Econometrica} \textbf{2000}, \emph{68}, 1097--1126.
\end{thebibliography}

\end{document}
